% GENERATED FILE. Edit canonical lessons, then run npm run generate:books.
% canonical-source-hash: fnv1a64:c987f631e9195503
% canonical-lessons: LA-C179-praesertim, LA-C179-potius, LA-C179-maxime, LA-C179-valde, LA-C179-admodum

\chapter{Especially, Rather, Most, Very, Quite}
\label{ch:la-especially-rather-most-very-quite}
\hlchaptermodality{179}

\begin{chapteropening}
\textbf{By the end of this chapter:} \emph{I can say especially, rather, most, very and quite.}

Complete the last lesson of chapter 179: I can say especially, rather, most, very and quite.
\end{chapteropening}

\section[praesertim]{praesertim — especially}

\label{lesson:LA-C179-praesertim}



\begin{quote}
Before the new one: say the Latin for even if, then the Latin for unless.
\end{quote}

\subsection*{You'll want to know: praesertim}
\textbf{praesertim} — \textquotedblleft{}especially\textquotedblright{}.

\begin{cousinweb}
\textbf{praesertim} is built on \emph{praeserō}, \textquotedblleft{}to put in front of a row\textquotedblright{} (Lewis \& Short), from \emph{prae}, \textquotedblleft{}before\textquotedblright{}, and \emph{serō}, \textquotedblleft{}to join, to line up\textquotedblright{}, with \emph{-tim}, an adverb ending: placed first, so especially.
\end{cousinweb}

\begin{grammarlens}[title={What you've built}]
One more word for this chapter.
\end{grammarlens}

\subsection*{Your turn}
\begin{itemize}
\raggedright
  \item \emph{Say it:} \emph{praesertim}
  \item \emph{Say it:} \emph{praesertim}, once more
  \item \emph{Say it:} say \emph{praesertim}
  \item \emph{Recall:} say the Latin for even if, then the Latin for unless, then say \emph{praesertim} again
\end{itemize}

\begin{tcolorbox}[breakable,colback=teal!4,colframe=teal!35!black,title={Before you move on}]
What does praesertim mean? (\textquotedblleft{}Especially\textquotedblright{}.) Say it once more.
\end{tcolorbox}
\section[potius]{potius — rather}

\label{lesson:LA-C179-potius}



\begin{quote}
Before the new one: say the Latin for unless, then the Latin for especially.
\end{quote}

\subsection*{You'll want to know: potius}
\textbf{potius} — \textquotedblleft{}rather\textquotedblright{}.

\begin{cousinweb}
\textbf{potius} is the comparative adverb of \emph{potis}, \textquotedblleft{}able, powerful\textquotedblright{} (Lewis \& Short): more powerfully, so rather.
\end{cousinweb}

\begin{grammarlens}[title={What you've built}]
One more word for this chapter.
\end{grammarlens}

\subsection*{Your turn}
\begin{itemize}
\raggedright
  \item \emph{Say it:} \emph{potius}
  \item \emph{Say it:} \emph{potius}, once more
  \item \emph{Say it:} say \emph{potius}
  \item \emph{Recall:} say the Latin for unless, then the Latin for especially, then say \emph{potius} again
\end{itemize}

\begin{tcolorbox}[breakable,colback=teal!4,colframe=teal!35!black,title={Before you move on}]
What does potius mean? (\textquotedblleft{}Rather\textquotedblright{}.) Say it once more.
\end{tcolorbox}
\section[maximē]{maximē — most, especially}

\label{lesson:LA-C179-maxime}



\begin{quote}
Before the new one: say the Latin for especially, then the Latin for rather.
\end{quote}

\subsection*{You'll want to know: maximē}
\textbf{maximē} — \textquotedblleft{}most, especially\textquotedblright{}.

\begin{cousinweb}
\textbf{maximē} is the adverb of the superlative of \emph{magnus}, \textquotedblleft{}big\textquotedblright{} (Lewis \& Short): \emph{maximus}, \textquotedblleft{}greatest\textquotedblright{}.
\end{cousinweb}

\begin{grammarlens}[title={What you've built}]
One more word for this chapter.
\end{grammarlens}

\subsection*{Your turn}
\begin{itemize}
\raggedright
  \item \emph{Say it:} \emph{maximē}
  \item \emph{Say it:} \emph{maximē}, once more
  \item \emph{Say it:} say \emph{maximē}
  \item \emph{Recall:} say the Latin for especially, then the Latin for rather, then say \emph{maximē} again
\end{itemize}

\begin{tcolorbox}[breakable,colback=teal!4,colframe=teal!35!black,title={Before you move on}]
What does maximē mean? (\textquotedblleft{}Most, especially\textquotedblright{}.) Say it once more.
\end{tcolorbox}
\section[valdē]{valdē — very}

\label{lesson:LA-C179-valde}



\begin{quote}
Before the new one: say the Latin for rather, then the Latin for most.
\end{quote}

\subsection*{You'll want to know: valdē}
\textbf{valdē} — \textquotedblleft{}very\textquotedblright{}.

\begin{cousinweb}
\textbf{valdē} is a shortened adverb of \emph{validus}, \textquotedblleft{}strong\textquotedblright{} (Lewis \& Short): strongly, so very.
\end{cousinweb}

\begin{grammarlens}[title={What you've built}]
One more word for this chapter.
\end{grammarlens}

\subsection*{Your turn}
\begin{itemize}
\raggedright
  \item \emph{Say it:} \emph{valdē}
  \item \emph{Say it:} \emph{valdē}, once more
  \item \emph{Say it:} say \emph{valdē}
  \item \emph{Recall:} say the Latin for rather, then the Latin for most, then say \emph{valdē} again
\end{itemize}

\begin{tcolorbox}[breakable,colback=teal!4,colframe=teal!35!black,title={Before you move on}]
What does valdē mean? (\textquotedblleft{}Very\textquotedblright{}.) Say it once more.
\end{tcolorbox}
\section[admodum]{admodum — quite, very}

\label{lesson:LA-C179-admodum}



\begin{quote}
Before the new one: say the Latin for most, then the Latin for very.
\end{quote}

\subsection*{You'll want to know: admodum}
\textbf{admodum} — \textquotedblleft{}quite, very\textquotedblright{}.

\begin{cousinweb}
\textbf{admodum} is \emph{ad}, \textquotedblleft{}up to\textquotedblright{}, with \emph{modum}, \textquotedblleft{}the measure\textquotedblright{} (Lewis \& Short): up to the full measure, so quite.
\end{cousinweb}

\begin{grammarlens}[title={What you've built}]
One more word for this chapter.
\end{grammarlens}

\subsection*{Your turn}
\begin{itemize}
\raggedright
  \item \emph{Say it:} \emph{admodum}
  \item \emph{Say it:} \emph{admodum}, once more
  \item \emph{Say it:} say \emph{admodum}
  \item \emph{Recall:} say the Latin for most, then the Latin for very, then say \emph{admodum} again
\end{itemize}

\begin{tcolorbox}[breakable,colback=teal!4,colframe=teal!35!black,title={Before you move on}]
What does admodum mean? (\textquotedblleft{}Quite, very\textquotedblright{}.) Say it once more.
\end{tcolorbox}
